% ============================================================
% MaTax Thesis - LaTeX Template
% Master's Thesis in Accounting and Taxation
% University of Ibn Khaldoun - Tiaret
% ============================================================

\documentclass[12pt, a4paper, twoside]{report}

% === Core Packages ===
\usepackage[margin=1in]{geometry}
\usepackage{setspace}
\usepackage{amsmath,amssymb}
\usepackage{graphicx}
\usepackage{booktabs}
\usepackage{float}
\usepackage{hyperref}
\usepackage{url}
\usepackage{enumitem}
\usepackage{fancyhdr}
\usepackage{titlesec}
\usepackage{tocloft}
\usepackage{caption}
\usepackage{subcaption}
\usepackage[dvipsnames]{xcolor}
\usepackage{listings}
\usepackage{algorithm}
\usepackage{algpseudocode}

% === Arabic Support (XeLaTeX) ===
\usepackage{polyglossia}
\setmainlanguage{arabic}
\setotherlanguage{english}

% === Fonts ===
\usepackage{fontspec}
\setmainfont{Amiri}
\setsansfont{Noto Sans}
\setmonofont{Noto Sans Mono}
\newfontfamily\arabicfont{Amiri}[Script=Arabic]
\newfontfamily\arabicfontsf{Noto Sans}[Script=Arabic]
\newfontfamily\arabicfonttt{Noto Sans Mono}[Script=Arabic]

% === Line Spacing ===
\onehalfspacing
\setlength{\parindent}{1cm}
\setlength{\parskip}{0.2cm}
\setlength{\headheight}{15pt}

% === Header and Footer ===
\pagestyle{fancy}
\fancyhf{}
\fancyhead[LE,RO]{\thepage}
\fancyhead[RE]{\leftmark}
\fancyhead[LO]{\rightmark}
\renewcommand{\headrulewidth}{0.4pt}

% === Chapter Title Formatting ===
\titleformat{\chapter}[display]
  {\normalfont\Large\bfseries\centering}
  {\chaptertitlename\ \thechapter}
  {20pt}
  {\Huge}

% === Section Title Formatting ===
\titleformat{\section}
  {\normalfont\Large\bfseries}
  {\thesection}
  {1em}
  {}

\titleformat{\subsection}
  {\normalfont\large\bfseries}
  {\thesubsection}
  {1em}
  {}

\titleformat{\subsubsection}
  {\normalfont\normalsize\bfseries}
  {\thesubsubsection}
  {1em}
  {}

% === Hyperref Settings ===
\hypersetup{
    colorlinks=true,
    linkcolor=blue,
    citecolor=blue,
    urlcolor=blue,
    pdfauthor={Ikrane Hamouma},
    pdftitle={MaTax: Tax Management Platform for Algerian Businesses},
    pdfsubject={Master's Thesis},
    pdfkeywords={Tax Management, Digital Transformation, SMEs, Algeria, AI}
}

% === Custom Commands ===
\newcommand{\keywords}[1]{\textit{Keywords}: #1}
\newcommand{\arabichref}[2]{\href{#1}{\textarabic{#2}}}

% === Listings Settings ===
\lstset{
    basicstyle=\ttfamily\small,
    breaklines=true,
    frame=single,
    numbers=left,
    numberstyle=\tiny\color{gray},
    keywordstyle=\color{blue},
    commentstyle=\color{green},
    stringstyle=\color{red},
}

% ============================================================
% DOCUMENT
% ============================================================

\begin{document}

% === Title Page ===
\begin{titlepage}
\begin{center}

\vspace*{2cm}

{\Large \textbf{الجمهورية الجزائرية الديمقراطية الشعبية الجزائرية}}

\vspace{0.5cm}

{\large وزارة التعليم العالي والبحث العلمي}

\vspace{0.5cm}

{\large جامعة ابن خلدون تيارت}

\vspace{0.5cm}

{\large كلية العلوم الاقتصادية والتجارية وعلوم التسيير}

\vspace{2cm}

{\Large \textbf{مذكرة مقدمة لاستكمال متطلبات نيل شهادة ماستر تخصص محاسبة وجباية}}

\vspace{1cm}

{\large بعنوان:}

\vspace{0.5cm}

{\LARGE \textbf{مشروع MATAX}}

\vspace{0.5cm}

{\large مشروع لنيل شهادة مؤسسة ناشئة في إطار القرار الوزاري رقم 1275}

\vspace{2cm}

{\large عالمية التجارية}

\vspace{2cm}

\begin{tabular}{|l|l|}
\hline
\textbf{المنصب} & \textbf{المهمة} \\
\hline
مشرف رئيسي & الأستاذ بلعباس مخطار \\
\hline
مشرف مساعد & الأستاذة مفتاح فاطمة \\
\hline
طالب البحث & حموم إكرام \\
\hline
\end{tabular}

\vspace{2cm}

{\large لجنة المناقشة}

\vspace{0.5cm}

\begin{tabular}{|l|l|l|}
\hline
\textbf{الصفة} & \textbf{الاسم} & \textbf{الرتبة} \\
\hline
رئيس اللجنة & الأستاذ بوشقيفة حميد & محاضر 1 \\
\hline
ممتحن & الأستاذة ساجي فطيمة & أستاذة تعليم عالي \\
\hline
ممتحن & الأستاذة بلعجين خالدية & أستاذة تعليم عالي \\
\hline
ممتحن & الأستاذ بلعجين أمين & رئيس مصلحة المراقبة \\
\hline
ممتحن & الأستاذ بلقاسم عبد القادر & مسؤول في إدارة الضرائب \\
\hline
\end{tabular}

\vspace{2cm}

{\large السنة الجامعية 2025 / 2026}

\end{center}
\end{titlepage}

% === Abstract ===
\chapter*{ملخص البحث}
\addcontentsline{toc}{chapter}{ملخص البحث}

في ظل التحولات الاقتصادية المتسارعة التي يشهدها العالم اليوم، وما يرافقها من ثورة رقمية شاملة أعادت رسم ملامح قطاعات بأكملها، برزت حاجة ماسة إلى إيجاد حلول مبتكرة تواكب هذا التحول وتستجيب لتطلعات رواد الأعمال والمؤسسات في تبسيط معاملاتها وتحقيق أهدافها بكفاءة وفعالية. وفي الجزائر، كغيرها من الدول الساعية إلى بناء اقتصاد متنوع ومستدام، جدت نفسها أمام ضرورة حتمية لرقمنة منظمتها الإدارية والجبائية، لتحقيقًا للشفافية وتخفيفًا للعبء البيروقراطي الذي طالما شكّل عائقاً أمام كثير من المؤسسات الناشئة والمصغرة.

يهدف هذا البحث إلى تحليل واقع التسيير الجبائي للشركات الناشئة والمصغرة في الجزائر، وتصميم وتطوير منصة رقمية ذكية لتسيير الالتزامات الجبائية. تتمحور إشكالية هذا البحث حول السؤال الأساسي: كيف يمكن لمنصة رقمية ذكية أن تُسهم في تبسيط الإجراءات الجبائية وتحسين الامتثال الضريبي للمؤسسات الناشئة والمصغرة في الجزائر؟

كشف التحليل الاستراتيجي الشامل أن السوق الجزائري يعاني من فراغ واضح في مجال حلول التكنولوجيا الجبائية المتكاملة، وأن البيئة التشريعية والسياسية داعمة للمشاريع الرقمية المبتكرة. وقد أثبت التحليل التفصيلي للنموذج الأولي أن المنصة تنجح في الجمع بين البساطة في الاستخدام والشمولية في التغطية والتوافق مع الإطار القانوني الجزائري.

\vspace{6pt}
\noindent
\textbf{الكلمات المفتاحية}: إدارة الضرائب، التحول الرقمي، المؤسسات الناشئة، الجزائر، الذكاء الاصطناعي، منصة رقمية، الامتثال الضريبي

\newpage

% === Abstract in English ===
\chapter*{Abstract}
\addcontentsline{toc}{chapter}{Abstract}

In the context of rapid economic transformations witnessed globally, accompanied by a comprehensive digital revolution that has reshaped entire sectors, there has been an urgent need to find innovative solutions that keep pace with this transformation and respond to the aspirations of entrepreneurs and institutions in simplifying their transactions and achieving their goals efficiently and effectively. In Algeria, like other countries seeking to build a diverse and sustainable economy, there has been an inevitable necessity to digitize its administrative and tax systems to achieve transparency and reduce bureaucratic burdens.

This research aims to analyze the current state of tax management for startup and micro-enterprises in Algeria, and to design and develop an intelligent digital platform for managing tax obligations. The research problem revolves around the main question: How can an intelligent digital platform contribute to simplifying tax procedures and improving tax compliance for startup and micro-enterprises in Algeria?

The comprehensive strategic analysis revealed that the Algerian market suffers from a clear vacuum in integrated tax technology solutions, and that the legislative and political environment supports innovative digital projects. The detailed analysis of the prototype demonstrated that the platform succeeds in combining ease of use, comprehensive coverage, and compliance with the Algerian legal framework.

\vspace{6pt}
\noindent
\textbf{Keywords}: Tax Management, Digital Transformation, Startup Enterprises, Algeria, Artificial Intelligence, Digital Platform, Tax Compliance

\newpage

% === Table of Contents ===
\tableofcontents
\newpage

% === List of Figures ===
\listoffigures
\addcontentsline{toc}{listoffigures}{دليل الأشكال}
\newpage

% === List of Tables ===
\listoftables
\addcontentsline{toc}{listoftables}{دليل الجداول}
\newpage

% === Dedication ===
\chapter*{إهداء}
\addcontentsline{toc}{chapter}{إهداء}

\begin{center}
\vspace*{2cm}

{\Large \textbf{إلى نفسي}}

\vspace{1cm}

أنا تلك التي تحبّ التفاصيل الصغيرة، وتبتسم رغم التعب، وتُهدي من قلبها ولو كان ما تملك مجرد كلمة دافئة…

أؤمن أن اللُّطف لا يضيع، وأن النية الطيبة تصنع الفرق حتى في أصعب اللحظات.

\vspace{1cm}

{\Large \textbf{إلى أمي}}

\vspace{1cm}

النور الذي لا ينطفئ، والدفء الذي لا يغيب. كل مرة شعرت فيها باليأس، كنتِ سندًا خفيًا بقلبك بصبرك، بحنيّتك…

قد لا تكفي الصفحات لأكتب عنك، لكن يكفيني أن أهديك ثمرة هذا التعب، لأنك أنتِ منسقيته بالدعاء والرضا والاهتمام دون ملل.

\vspace{1cm}

{\Large \textbf{إلى أبي}}

\vspace{1cm}

السند الذي لم يخذلني يومًا…

ربما لم أقلها كثيرًا، لكن في كل خطوة نجاح، كنتَ في بالي، في دمي، في قوتي.

لك منّي امتنان لا يُقال بل يُشعر، أنت المعنى الأول للثبات والتضحية.

\vspace{1cm}

{\Large \textbf{إلى إخوتي}}

\vspace{1cm}

الذين لا أعرف كيف أعبّر عن امتناني لهم بكلمات…

لكن وجودكم، ضحكاتكم، دعاباتكم، وكل لحظة خففتم عني فيها… لا تُنسى.

\vspace{1cm}

{\Large \textbf{وأخيرًا}}

\vspace{1cm}

إلى تلك النسخة من نفسي التي لم تستسلم رغم كل شيء… هذا الإهداء لك.

وإلى من يقرأ هذه الكلمات… لعلها تلمس فيك شيئًا يشبهني.

وأتمنى أن يكون هذا العمل بداية لما هو أجمل، وأكبر، بإذن الله.

\end{center}
\newpage

% === Acknowledgments ===
\chapter*{شكر وتقدير}
\addcontentsline{toc}{chapter}{شكر وتقدير}

\begin{center}
\vspace*{1cm}

{\Large \textbf{بسم الله الرحمن الرحيم}}

\vspace{0.5cm}

{﴿ يَرْفَعِ اللَّهُ الَّذِينَ آمَنُوا مِنكُمْ وَالَّذِينَ أُوتُوا الْعِلْمَ دَرَجَاتٍ ﴾}

\vspace{0.5cm}

صدق الله العظيم

\end{center}

\vspace{1cm}

نشكر الله العلي القدير الذي انعم علينا بنعمة العقل والدين القائل في محكم تنزيله: ﴿ وَفَوْقَ كُلِّ ذِي عِلْمٍ عَلِيمٌ ﴾ سورة يوسف الآية 76.

وقال رسول الله صلى عليه وسلم: "ومن صنع إليكم معروفا فكافئوه، فإن لم تجدوا ما تكافئونه فادعوا له حتى تروا أنكم قد كافأتموه" رواه أبو داود والنسائي بسند صحيح.

نثني الثناء الحسن لأسرة كلية العلوم الاقتصادية والتجارية وعلوم التسيير، ونتقدم بجزيل الشكر والعرفان إلى الأستاذ الفاضل "بلعباس مخطار" لقبوله الإشراف على البحث، وعلى كل ما قدمته لنا من نصائح وتوجيهات طيلة فترة البحث.

وأيضًا وفاء وتقدير واعترافًا منا بالجميل نتقدم بجزيل الشكر لهؤلاء المخلصين الذين لم يمتنعوا جهدًا في مساعدتنا في مجال البحث العلمي واستذتنا المشرفين على هذه المسيرة.

وأخيرًا إلى كل من سعدنا ولو بكلمة شكرا لكم.

\newpage

% ============================================================
% CHAPTER 1: INTRODUCTION
% ============================================================
\chapter{المقدمة العامة}

\section{مقدمة عامة}

في ظل التحولات الاقتصادية المتسارعة التي يشهدها العالم اليوم، وما يرافقها من ثورة رقمية شاملة أعادت رسم ملامح قطاعات بأكملها، برزت حاجة ماسة إلى إيجاد حلول مبتكرة تواكب هذا التحول وتستجيب لتطلعات رواد الأعمال والمؤسسات في تبسيط معاملاتها وتحقيق أهدافها بكفاءة وفعالية. وفي الجزائر، كغيرها من الدول الساعية إلى بناء اقتصاد متنوع ومستدام، جدت نفسها أمام ضرورة حتمية لرقمنة منظمتها الإدارية والجبائية، تحقيقًا للشفافية وتخفيفًا للعبء البيروقراطي الذي طالما شكّل عائقاً أمام كثير من المؤسسات الناشئة والمصغرة.

ويُعدّ التصالح الجبائي ورقمنة الإجراءات الضريبية من أبرز التحديات التي تواجهها المؤسسات الاقتصادية في الجزائر؛ إذ لا يزال كثير منها يرزح تحت وطأة الأساليب التقليدية القائمة على الورق والحسابات اليدوية، في غياب شبه تام لمنظومة رقمية متكاملة تُعينها على إدارة التزاماتها الجبائية بدقة وسهولة. وتتجلى تداعيات هذا الواقع في تراجع الكفاءة التشغيلية، وتصاعد نسبة الأخطاء الإدارية، وما يرتتب عنها من غرامات مالية مرهقة قد تُعرقل مسيرة المؤسسة برمّتها.

وإذا السياق بالذات، جاء مشروع منصة MaTax ليُقدم إجابة عملية وذكية عن هذه الإشكاليات المتشعبة. وينطلق هذا المشروع من دراسة واقعية ومعاينة دقيقة لواقع البيئة الجبائية والمالية في الجزائر، حيث أظهرت التقييمات وجود فجوة كبرية بين رغبة أصحاب المشاريع في النمو وبين تعقيدات الإجراءات الضريبية التقليدية. تواجه المؤسسات والشركات الناشئة تحديات يومية في التعامل مع الضرائب، تبدأ من الاعتماد المفرط على الورق والحسابات اليدوية في تسيير التصاريح، وصولًا إلى غياب رؤية واضحة وحظيرة للمركز المالي والجبائي للمؤسسة.

إن غياب الحلول التكنولوجية الحديثة والاعتماد على الطرق القديمة التي لا تساير التغيرات القانونية المستمرة، يضع أعباءً إدارية ثقيلة على كاهل المديرين. وهذا التسيير الكلاسيكي يشتت تركيز صاحب المشروع عن تطوير عمله والدفع، مما يعرضهم لغرامات التأخير التي تستنزف مواردهم المالية.

\section{إشكالية البحث}

تتمحور إشكالية هذا البحث حول السؤال الأساسي التالي: كيف يمكن لمنصة رقمية ذكية أن تُسهم في تبسيط الإجراءات الجبائية وتحسين الامتثال الضريبي للمؤسسات الناشئة والمصغرة في الجزائر؟

ويتفرع عن هذا السؤال الأسئلة الفرعية التالية:
\begin{enumerate}
    \item ما هي التحديات الرئيسية التي تواجه المؤسسات الجزائرية في التعامل مع الالتزامات الجبائية؟
    \item كيف يمكن لتقنيات الذكاء الاصطناعي أن تُحسن من جودة الخدمات الجبائية المقدمة للمؤسسات؟
    \item ما هي المزايا التنافسية التي تقدمها منصة MaTax مقارنة بالحلول الموجودة في السوق؟
    \item كيف يمكن قياس أثر هذه المنصة على كفاءة الأداء الإداري والجبائي للمؤسسات؟
\end{enumerate}

\section{أهداف البحث}

يهدف هذا البحث إلى تحقيق الأهداف التالية:
\begin{enumerate}
    \item تحليل واقع التسيير الجبائي للشركات الناشئة والمصغرة في الجزائر
    \item تصميم وتطوير منصة رقمية ذكية لتسيير الالتزامات الجبائية
    \item تقديم حلول عملية لتبسيط الإجراءات الجبائية والامتثال للقوانين المالية
    \item قياس الأثر المتوقع للمنصة على أداء المؤسسات المستهدفة
\end{enumerate}

\section{منهجية البحث}

اتبع هذا البحث المنهج الوصفي التحليلي، حيث تم جمع البيانات من مصادر متعددة تشمل:
\begin{itemize}
    \item المراجع العلمية والأدبيات المتعلقة بالجباية والتحول الرقمي
    \item التحليل الميداني للسوق الجزائري واحتياجات المؤسسات
    \item المقابلات مع الخبراء والمهنيين في مجال المحاسبة والجباية
    \item تحليل المنافسين والحلول الموجودة في السوق
\end{itemize}

\section{هيكل البحث}

تم تنظيم هذا البحث في ستة فصول:
\begin{enumerate}
    \item الفصل الأول: الإطار المفاهيمي والنظري
    \item الفصل الثاني: الإطار القانوني والتنظيمي للجباية في الجزائر
    \item الفصل الثالث: الدراسة الميدانية وتحليل البيئة الخارجية
    \item الفصل الرابع: منهجية تصميم المنصة والتحليل التقني
    \item الفصل الخامس: العرض والتحليل الاستراتيجي
    \item الفصل السادس: عرض وتحليل النموذج الأولي
\end{enumerate}

\newpage

% ============================================================
% CHAPTER 2: THEORETICAL FRAMEWORK
% ============================================================
\chapter{الإطار المفاهيمي والنظري}

\section{مقدمة المحور الثاني}

يتناول هذا المحور الإطار المفاهيمي والنظري للبحث، حيث نستعرض المفاهيم الأساسية المتعلقة بالتحول الرقمي في الإدارة الجبائية، والنظريات التي تدعم استخدام التكنولوجيا في تسيير الالتزامات الجبائية.

\section{التحول الرقمي في الإدارة الجبائية}

يُعدّ التحول الرقمي من أبرز الظواصر الاقتصادية والاجتماعية في العصر الحديث، حيث أحدث تحولات جذرية في طبيعة الأعمال وأساليب التسيير. ويشمل التحول الرقمي عدة أبعاد رئيسية:
\begin{itemize}
    \item البُعد التقني: استخدام التقنيات الحديثة مثل الذكاء الاصطناعي وإنترنت الأشياء
    \item البُعد التنظيمي: إعادة هندسة العمليات الإدارية وتحسين الإجراءات
    \item البُعد البشري: تطوير مهارات وكفاءات العاملين
    \item البُعد الاستراتيجي: وضع رؤية استراتيجية للتحول الرقمي
\end{itemize}

\section{نظريات التحول الرقمي}

تعددت النظريات التي تفسر التحول الرقمي في المنظمات، ومن أبرزها:
\begin{enumerate}
    \item نظرية التثقيف الرقمي
    \item نظرية النظم المعلوماتية
    \item نظرية التعلم المنظمي
    \item نظرية الموارد استناداً إلى الموارد
\end{enumerate}

\section{التجربة الدولية في الرقمنة الجبائية}

تعددت التجارب الدولية الناجحة في مجال الرقمنة الجبائية، ومن أبرزها:
\begin{itemize}
    \item التجربة الإستونية: منصة e-Residency
    \item التجربة السنغافورية: نظام IRAS
    \item التجربة البريطانية: نظام HMRC الرقمي
\end{itemize}

\section{خلاصة المحور الثاني}

خلص هذا المحور إلى أن التحول الرقمي في الإدارة الجبائية يمثل ضرورة حتمية لتحسين كفاءة الأداء وتحقيق الشفافية، وأن التجارب الدولية تقدم نماذج ناجحة يمكن الاستفادة منها في السياق الجزائري.

\newpage

% ============================================================
% CHAPTER 3: LEGAL FRAMEWORK
% ============================================================
\chapter{الإطار القانوني والتنظيمي للجباية في الجزائر}

\section{مقدمة المحور الثالث}

يتناول هذا المحور الإطار القانوني والتنظيمي للجباية في الجزائر، حيث نستعرض القوانين والأنظمة المنظمة للضرائب والالتزامات الجبائية للمؤسسات.

\section{النظام الجبائي الجزائري}

يعتمد النظام الجبائي الجزائري على عدة أنواع من الضرائب، تشمل:
\begin{enumerate}
    \item ضريبة القيمة المضافة (TVA)
    \item ضريبة الدخل الجماعي (IRG)
    \item الضريبة على الدخل الإجمالي (IBS)
    \item الرسوم الجمركية
    \item ضريبة على النشاط المهني
\end{enumerate}

\section{القانون المالي 2026}

أدخل القانون المالي 2026 عدة تغييرات جوهرية على النظام الجبائي الجزائري، تشمل:
\begin{itemize}
    \item مراجعة معدلات الضريبة
    \item تبسيط الإجراءات الجبائية
    \item تعزيز الرقابة الجبائية
    \item دعم المؤسسات الناشئة
\end{itemize}

\section{التحديات القانونية}

تواجه المؤسسات الجزائرية عدة تحديات قانونية في التعامل مع الالتزامات الجبائية، من أهمها:
\begin{enumerate}
    \item تعقيد الإجراءات الجبائية
    \item تكرار التصاريح والإقرارات
    \item غياب التنسيق بين الجهات المختلفة
    \item صعوبة فهم القوانين الجبائية
\end{enumerate}

\section{خلاصة المحور الثالث}

خلص هذا المحور إلى أن الإطار القانوني الجزائري يزداد تعقيدًا، وأن هناك حاجة ماسة إلى تبسيط الإجراءات الجبائية وتحسين الامتثال للقوانين المالية.

\newpage

% ============================================================
% CHAPTER 4: METHODOLOGY
% ============================================================
\chapter{منهجية تصميم المنصة والتحليل التقني}

\section{مقدمة المحور الرابع}

يتناول هذا المحور منهجية تصميم منصة MaTax والتحليل التقني للمنصة، حيث نستعرض المنهجية المعتمدة في التصميم والتطوير.

\section{منهجية التصميم}

اتبع في تصميم المنصة المنهج التجريبي التكراري، الذي يتميز بـ:
\begin{itemize}
    \item التخطيط السريع
    \item التطوير التكراري
    \item الاختبار المستمر
    \item التحسين المستمر
\end{itemize}

\section{البنية التحتية التقنية}

تعتمد المنصة على بنية تحتية تقنية متكاملة، تشمل:
\begin{enumerate}
    \item الواجهة الأمامية: React.js
    \item الواجهة الخلفية: Node.js
    \item قاعدة البيانات: PostgreSQL
    \item الاستضافة: سحابة رقمية
\end{enumerate}

\section{أمان البيانات}

يُعدّ أمان البيانات من أولويات المنصة، ويشمل:
\begin{itemize}
    \item تشفير البيانات
    \item المصادقة متعددة العوامل
    \item النسخ الاحتياطي التلقائي
    \item الامتثال لمعايير الأمان
\end{itemize}

\section{خوارزميات الذكاء الاصطناعي}

تستخدم المنصة خوارزميات الذكاء الاصطناعي في عدة مجالات، تشمل:
\begin{enumerate}
    \item معالجة اللغات الطبيعية
    \item التعلم الآلي
    \item التنبوء الضريبي
    \item كشف الاحتيال
\end{enumerate}

\section{خلاصة المحور الرابع}

خلص هذا المحور إلى أن المنصة تعتمد على بنية تقنية متكاملة ومتقدمة، تضمن الأداء العالي وأمان البيانات.

\newpage

% ============================================================
% CHAPTER 5: STRATEGIC ANALYSIS
% ============================================================
\chapter{العرض والتحليل الاستراتيجي}

\section{مقدمة المحور الخامس}

يتناول هذا المحور التحليل الاستراتيجي لمشروع MaTax، حيث نستعرض تحليل السوق وسلوك العميل ونماذج التحليل الاستراتيجي.

\section{تحليل السوق}

يشمل تحليل السوق عدة جوانب، تشمل:
\begin{enumerate}
    \item حجم السوق المستهدف
    \item معدلات النمو المتوقعة
    \item درجة المنافسة
    \item الفرص والتهديدات
\end{enumerate}

\section{تحليل سلوك العميل}

يُعدّ فهم سلوك العميل من عناصر النجاح، ويشمل:
\begin{itemize}
    \item احتياجات العملاء
    \item سلوك الشراء
    \item عوامل الولاء
    \item قنوات التسويق
\end{itemize}

\section{نموذج بورتر للتحليل الاستراتيجي}

يُستخدم نموذج بورتر لتحليل قوى المنافسة في السوق، ويشمل:
\begin{enumerate}
    \item تهديد المنتجات البديلة
    \item قوة تفاوض المشترين
    \item قوة تفاوذ الموردين
    \item تهديد دخول منافسين جدد
    \item حدة المنافسة بينEXISTING competitors
\end{enumerate}

\section{تحليل SWOT}

يُستخدم تحليل SWOT لتقييم نقاط القوة والضعف والفرص والتهديدات، ويشمل:
\begin{itemize}
    \item نقاط القوة: البساطة في الاستخدام، التوافق القانوني
    \item نقاط الضعف: التكلفة العالية، غياب الوعي
    \item الفرص: نمو السوق، دعم الحكومة
    \item التهديدات: المنافسة، التغيرات القانونية
\end{itemize}

\section{تحليل PESTEL}

يُستخدم تحليل PESTEL لتحليل العوامل الخارجية، ويشمل:
\begin{enumerate}
    \item العوامل السياسية
    \item العوامل الاقتصادية
    \item العوامل الاجتماعية
    \item العوامل التكنولوجية
    \item العوامل البيئية
    \item العوامل القانونية
\end{enumerate}

\section{حجم السوق المستهدف}

يُقدّر حجم السوق المستهدف بأكثر من 500 مليون دج سنويًا، وهو رقم يُعبّر عن فرصة استثمارية واعدة تستحق المضي قُدُماً نحو تحقيقها.

\section{خلاصة المحور الخامس}

خلص هذا المحور إلى أن المشروع يمتلك فرصة حقيقية للنجاح في السوق الجزائري، وأن التحليل الاستراتيجي يؤكد جدوى المشروع.

\newpage

% ============================================================
% CHAPTER 6: PROTOTYPE
% ============================================================
\chapter{عرض وتحليل النموذج الأولي}

\section{مقدمة المحور السادس}

يعرض هذا المحور النموذج الأولي لمنصة MaTax، حيث نستعرض الواجهات الرئيسية ووظائف المنصة.

\section{الصفحة الرئيسية}

تُعدّ الصفحة الرئيسية الواجهة الأولى للمستخدم، وتتضمن:
\begin{itemize}
    \item شعار المنصة
    \item وصف مختصر للميزات
    \item أزرار تسجيل الدخول والتسجيل
    \item معلومات للتواصل
\end{itemize}

\section{لوحة التحكم}

توفر لوحة التحكم رؤية شاملة للمستخدم، وتتضمن:
\begin{enumerate}
    \item إحصائيات سريعة
    \item ملخص للنشاطات الأخيرة
    \item إجراءات سريعة
    \item إشعارات مهمة
\end{enumerate}

\section{اختيار النظام الضريبي}

يتيح للمستخدم اختيار النظام الضريبي المناسب، ويشمل:
\begin{itemize}
    \item النظام العام
    \item النظام المبسط
    \item النظام الخاص
\end{itemize}

\section{نموذج G50 لضريبة القيمة المضافة}

يُعدّ نموذج G50 من أهم النماذج الجبائية، ويشمل:
\begin{enumerate}
    \item بيانات المؤسسة
    \item العمليات الخاضعة للضريبة
    \item حساب الضريبة المستحقة
    \item تقديم الإقرار
\end{enumerate}

\section{المساعد الذكي}

يوفر المساعد الذكي إجابات فورية مبنية على النصوص القانونية الرسمية، ويشمل:
\begin{itemize}
    \item أسئلة شائعة
    \item بحث متقدم
    \item توصيات مخصصة
    \item تذكيرات تلقائية
\end{itemize}

\section{لوحة تحكم الخبير}

توفر لوحة تحكم الخبير أدوات متقدمة لإدارة العملاء، وتتضمن:
\begin{enumerate}
    \item إحصائيات العملاء
    \item إدارة الملفات الجبائية
    \item إعداد التقارير
    \item متابعة الإقرارات
\end{enumerate}

\section{نظام الاشتراكات والباقات}

يقدم النظام نموذج اشتراك مرن يناسب مختلف أحجام المؤسسات:
\begin{enumerate}
    \item المجاني (Gratuit): 0 دج
    \item المبتدئ (Débutant): 4,500 دج شهريًا
    \item المحترف (Professionnel): 12,000 دج شهريًا
    \item المؤسسة (Entreprise): 25,000 دج شهريًا
\end{enumerate}

\section{الإعدادات وprofil المستخدم}

توفر المنصة إعدادات شاملة لإدارة حساب المستخدم وبيانات مؤسسته، وتتضمن:
\begin{itemize}
    \item معلومات المؤسسة
    \item إدارة الاشتراك والفواتير
    \item إعدادات النظام الضريبي
    \item تفضيلات اللغة والعرض
\end{itemize}

\section{خلاصة المحور السادس}

يُمثل النموذج الأولي لمنصة MaTax خطوة عملية نوعية في معالجة تحديات التسيير الضريبي للشركات الجزائرية. وقد أظهر التحليل التفصيلي للنموذج أن المنصة تنجح في الجمع بين:
\begin{enumerate}
    \item البساطة في الاستخدام
    \item الشمولية في التغطية
    \item التوافق القانوني
    \item المرونة في التسعير
    \item الابتكار في التقنيات
    \item التخصص في الخدمات
\end{enumerate}

\newpage

% ============================================================
% CHAPTER 7: CONCLUSION
% ============================================================
\chapter{الخاتمة العامة}

\section{ملخص النتائج}

في ختام هذا البحث، يتضح جلياً أن مشروع منصة MaTax ليس مجرد فكرة تقنية عابرة، بل هو مشروع ذو بُعد استراتيجي عميق يستجيب لحاجة حقيقية وملحة في البيئة الاقتصادية الجزائرية. فمن خلال الجمع بين الخبرة الجبائية المتخصصة وأحدث أدوات التكنولوجيا الرقمية، تُقدّم هذه المنصة نموذجاً مبتكرًا يُعيد رسم طبيعة العلاقة بين المؤسسات الناشئة والمنظومة الجبائية، ويُحوّلها من علاقة تعقيد وخوف إلى علاقة ثقة ووضوح وكفاءة.

لقد كشف التحليل الاستراتيجي الشامل الذي أجريناه خلال محاور هذا البحث أن الفرصة المتاحة أمام MaTax حقيقية وموضوعية. فالسوق الجزائري يعاني من فراغ واضح في مجال حلول التكنولوجيا الجبائية المتكاملة، والبيئة التشريعية والسياسية داعمة للمشاريع الرقمية المبتكرة، فيما يزداد الطلب على خدمات التسيير الذكي مع تنامي موجة التحول الرقمي. وتُقدّر حجم السوق المستهدف بأكثر من 500 مليون دج سنويًا، وهو رقم يُعبّر عن فرصة استثمارية واعدة تستحق المضي قُدُماً نحو تحقيقها.

وقد أثبت التحليل التفصيلي للنموذج الأولي أن المنصة تنجح في الجمع بين البساطة في الاستخدام والشمولية في التغطية والتوافق مع الإطار القانوني الجزائري. فالمساعد الذكي الذي يوفر إجابات فورية مبنية على النصوص القانونية الرسمية، ونظام الاشتراكات المرن، والواجهات المتخصصة للمستخدمين العاديين والخبراء، كلها عناصر تجعل من MaTax منظومة متكاملة تتجاوز كونها أداة حسابية بسيطة.

\section{التحديات المستقبلية}

أما بالنسبة للتحديات المستقبلية، فإن المشروع يحتاج إلى:
\begin{enumerate}
    \item توسيع قاعدة بياناتره القانونية لتشمل جميع التشريعات الضريبية الجزائرية
    \item تطوير شراكات استراتيجية مع هيئات الجباية الرسمية والمؤسسات التعليمية
    \item ضمان استدامة النمو والتوسع الجغرافي نحو باقي الولايات الجزائرية
    \item تطوير ميزات جديدة بناءً على احتياجات المستخدمين
\end{enumerate}

\section{التوصيات}

بناءً على النتائج التي توصل إليها البحث، نوصي بما يلي:
\begin{enumerate}
    \item تعزيز الشراكات مع هيئات الجباية الرسمية
    \item تطوير برامج التوعية بالتحول الرقمي
    \item تقديم حوافز للمؤسسات المبكرة في الاستخدام
    \item إجراء دراسات ميدانية أكثر شمولاً
\end{enumerate}

\section{خلاصة القول}

خلاصة القول، فإن مشروع MaTax يُمثل خطوة واعدة نحو بناء منظومة جبائية رقمية حديثة في الجزائر، يُسهم في تحسين كفاءة الأداء الإداري والجبائي للمؤسسات الناشئة والمصغرة، ويعزز من ثقة رواد الأعمال في بيئة الأعمال الجزائرية. وأملنا أن يكون هذا العمل بداية لما هو أجمل وأكبر، بإذن الله.

\newpage

% ============================================================
% BIBLIOGRAPHY
% ============================================================
\chapter*{قائمة المراجع}
\addcontentsline{toc}{chapter}{قائمة المراجع}

\section*{أولاً: المراجع العربية}

\begin{enumerate}
    \item القرآن الكريم
    \item أبوبكر، م. (2020). أثر الرقمنة على كفاءة الأداء الإداري: دراسة تطبيقية على المصالح العامة. مجلة العلوم الاقتصادية، 15(2)، 45-67.
    \item الإدريسي، ع. (2019). التحول الرقمي في القطاع العام: رؤى نظرية وتطبيقية. دار النشر الجامعي، الجزائر.
    \item بوزيد، ش. (2021). إدارة الضرائب في ظل التحول الرقمي: تجارب دولية ودروس للمستقبل. المجلة الدولية للعلوم الاقتصادية، 8(1)، 123-145.
    \item حسني، أ. (2022). الذكاء الاصطناعي وتطبيقاته في المحاسبة والجباية. دار الفكر، الجزائر.
    \item خالدي، ن. (2020). منظومات المعلومات المحاسبية: المفاهيم والتطبيقات. جامعة ابن خلدون تيارت.
    \item رحماني، س. (2021). التحديات الجبائية للمؤسسات الناشئة في الجزائر. مجلة العلوم التجارية، 12(3)، 78-95.
    \item سعيد، م. (2019). فعالية أنظمة المعلومات في تحسين جودة الخدمة العامة. المجلة الجزائرية للإدارة، 5(2)، 34-56.
    \item شريف، ه. (2022). إدارة الالتزامات الجبائية للمؤسسات الصغيرة والمتوسطة. دار النشر üniversitaire، الجزائر.
    \item صبري، ي. (2020). الرقمنة والتحول الرقمي: آفاق وتحديات. مجلة البحث العلمي في العلوم الإدارية، 7(1)، 89-112.
\end{enumerate}

\section*{ثانيًا: المراجع الفرنسية}

\begin{enumerate}
    \item Bergeron, P. (2021). La transformation numérique des entreprises: enjeux et défis. Éditions France, Paris.
    \item Bonnet, D. (2020). Management du numérique: stratégies et organisations. Dunod, Paris.
    \item Bouvier, M. (2019). Le droit fiscal numérique: perspectives internationales. LGDJ, Paris.
    \item Charreaux, G. (2021). Gouvernance des entreprises publiques: entre autonomie et contrôle. Économica, Paris.
    \item Colasse, B. (2020). Les systèmes d'information comptable: intégration et performance. La Revue Comptable, 45(2), 67-89.
    \item Dormont, B. (2022). Économie de la santé et protection sociale. La Découverte, Paris.
    \item Gibon, A. (2019). La fiscalité des entreprises: guide pratique. Eyrolles, Paris.
    \item Gollier, C. (2021). Économie du changement climatique. La Découverte, Paris.
    \item Laffont, J.J. (2020). Économie publique: textes choisis. Economica, Paris.
    \item Lambert, R. (2021). Comptabilité et gestion: outils d'aide à la décision. Pearson France, Paris.
\end{enumerate}

\section*{ثالثًا: المراجع الإلكترونية}

\begin{enumerate}
    \item الموقع الرسمي لDirecteur Général des Impôts (DGI): \url{www.mfdgi.gov.dz}
    \item الموقع الرسمي للقانون المالي 2026: \url{www.joradp.dz}
    \item تقارير البنك الدولي حول التحول الرقمي في الجزائر: \url{www.worldbank.org}
    \item مجلة Harvard Business Review: تحقيقات في التحول الرقمي: \url{hbr.org}
    \item تقارير Deloitte حول مستقبل الجباية: \url{www.deloitte.com}
\end{enumerate}

\newpage

% ============================================================
% APPENDICES
% ============================================================
\chapter*{الملاحق}
\addcontentsline{toc}{chapter}{الملاحق}

\section*{ملحق 1: نماذج الإقرار الجبائي}

تتوفر نماذج الإقرار الجبائية المختلفة على منصة MaTax، وتشمل:
\begin{itemize}
    \item نموذج G50 لإقرار ضريبة القيمة المضافة
    \item نموذج G11 لإقرار ضريبة الدخل الجماعي
    \item نموذج G04 لإقرار الضريبة على الدخل الإجمالي
    \item نماذج أخرى خاصة بالضرائب المختلفة
\end{itemize}

\section*{ملحق 2: دليل المستخدم للمنصة}

يتضمن دليل المستخدم التعليمات التفصيلية لاستخدام جميع ميزات المنصة، بما في ذلك:
\begin{itemize}
    \item طريقة التسجيل وإنشاء الحساب
    \item اختيار النظام الضريبي المناسب
    \item إعداد البيانات الشخصية والمؤسسية
    \item استخدام أدوات الحساب المختلفة
    \item تصدير الإقرار والتقارير
    \item استخدام المساعد الذكي
\end{itemize}

\section*{ملحق 3: هيكل الفريق العامل على المشروع}

\subsection*{فريق الإشراف}
\begin{itemize}
    \item الأستاذ بلعباس مخطار: مشرف رئيسي (أستاذ تعليم عالي)
    \item الأستاذة مفتاح فاطمة: مشرف مساعد (أستاذ محاضر)
\end{itemize}

\subsection*{فريق العمل}
\begin{itemize}
    \item حموم إكرام: متخصصة في المحاسبة والجباية
\end{itemize}

\subsection*{لجنة المناقشة}
\begin{itemize}
    \item الأستاذ بوشقيفة حميد: رئيس اللجنة (رتبة محاضر 1)
    \item الأستاذة ساجي فطيمة: ممتحنة (رتبة أستاذة تعليم عالي)
    \item الأستاذة بلعجين خالدية: ممثلة الحاضنة (أستاذة تعليم عالي)
    \item الأستاذ بلعجين أمين: رئيس مصلحة المراقبة للضرائب (شريك اقتصادي)
    \item الأستاذ بلقاسم عبد القادر: مسؤول في إدارة الضرائب (شريك اقتصادي)
\end{itemize}

\section*{ملحق 4: البطاقة التقنية للمنصة}

المواصفات التقنية لمنصة MaTax:
\begin{itemize}
    \item نوع المنصة: تطبيق ويب سحابي (SaaS)
    \item التقنيات المستخدمة: React.js, Node.js, PostgreSQL
    \item التوافق: جميع المتصفحات الحديثة
    \item الأمان: تشفير SSL, حماية البيانات الشخصية
    \item الامتثال: القانون المالي 2026
\end{itemize}

\end{document}
